% अभ्यास-संग्रह — अध्याय 4
% अध्याय इस फ़ाइल के नाम से निर्धारित होता है।

\begin{exo}[किसी पथ के अनुदिश अवकल रूप का समाकलन]{integrale-forme-differentielle}
\keywords{अवकल रूप, यथार्थ अवकल, श्वार्ज कसौटी, रेखा समाकल}

दो अवकल रूप
\[
\omega_1 = y\,\mathrm{d}x + x\,\mathrm{d}y,
\qquad
\omega_2 = y\,\mathrm{d}x,
\]
और मूलबिंदु $O(0,0)$ को $M(1,1)$ से जोड़ने वाले तीन पथ लें: $\gamma_1$ में
क्षैतिज खंड $O\to(1,0)$ के बाद ऊर्ध्वाधर खंड $(1,0)\to M$ है; $\gamma_2$ में
ऊर्ध्वाधर खंड $O\to(0,1)$ के बाद क्षैतिज खंड $(0,1)\to M$ है; और $\gamma_3$
$O$ तथा $M$ को सीधे जोड़ने वाला रेखाखंड है।

\begin{enumerate}
\item ऐसा फलन $f$ देकर दिखाएँ कि $\omega_1 = \mathrm{d}f$ है, अतः $\omega_1$ यथार्थ है।
\item प्रत्येक खंड का प्राचलीकरण करके $\int_{\gamma_1}\omega_1$ और $\int_{\gamma_2}\omega_1$ निकालें। बिना गणना भी परिणाम पाएँ।
\item $\omega_2$ पर श्वार्ज कसौटी लगाएँ। क्या निष्कर्ष मिलता है?
\item $\int_{\gamma_1}\omega_2$, $\int_{\gamma_2}\omega_2$ और $\int_{\gamma_3}\omega_2$ निकालकर टिप्पणी करें।
\end{enumerate}

\begin{indication}
क्षैतिज खंड पर $\mathrm{d}y = 0$ और ऊर्ध्वाधर खंड पर $\mathrm{d}x = 0$ होता है;
हर समाकल साधारण एकचर समाकल में बदल जाता है।
\end{indication}

\begin{solution}
\textbf{प्रश्न 1।} $f(x,y) = xy$ उपयुक्त है, क्योंकि $\partial f/\partial x = y$ तथा $\partial f/\partial y = x$।

\textbf{प्रश्न 2।} $t\mapsto\bigl(x(t),y(t)\bigr)$ से प्राचलित वक्र के लिए
\[
\int_\gamma\omega_1
= \int \left[y(t)x'(t)+x(t)y'(t)\right],\mathrm{d}t.
\]
$\gamma_1$ के पहले खंड पर $x(t)=t$, $y(t)=0$ और दूसरे पर $x(t)=1$, $y(t)=t$ है;
दोनों में $t\in[0,1]$। अतः
\[
\begin{aligned}
I_1=\int_{\gamma_1}\omega_1
&= \int_0^1(0\times1+t\times0)\,\mathrm{d}t
+\int_0^1(t\times0+1\times1)\,\mathrm{d}t\\
&=0+\int_0^1\mathrm{d}t=1.
\end{aligned}
\]
$\gamma_2$ के पहले खंड पर $x(t)=0$, $y(t)=t$ और दूसरे पर $x(t)=t$, $y(t)=1$ है। इसलिए
\[
\begin{aligned}
I_2=\int_{\gamma_2}\omega_1
&= \int_0^1(t\times0+0\times1)\,\mathrm{d}t
+\int_0^1(1\times1+t\times0)\,\mathrm{d}t\\
&=0+\int_0^1\mathrm{d}t=1.
\end{aligned}
\]
दोनों मान समान हैं। $\omega_1=\mathrm{d}f$ और $f(x,y)=xy$ होने से समाकल केवल सिरों पर निर्भर है:
\[
\int_\gamma\omega_1=f(M)-f(O)=f(1,1)-f(0,0)=1.
\]

\textbf{प्रश्न 3।} $\omega_2=A(x,y)\,\mathrm{d}x+B(x,y)\,\mathrm{d}y$ लिखें। यहाँ
$A(x,y)=y$, $B(x,y)=0$। यथार्थ रूप के लिए श्वार्ज कसौटी देती
\[
\frac{\partial A}{\partial y}=\frac{\partial B}{\partial x}.
\]
परंतु $\partial A/\partial y=1$ और $\partial B/\partial x=0$ हैं। मिश्रित अवकलज असमान हैं,
अतः $\omega_2$ यथार्थ नहीं है और उसका समाकल $O$ से $M$ तक चुने गए पथ पर निर्भर हो सकता है।

\textbf{प्रश्न 4।} $\omega_2=y\,\mathrm{d}x$ में केवल अशून्य $y$ पर $x$ का परिवर्तन योगदान देता है। अतः
\[
\begin{aligned}
J_1=\int_{\gamma_1}\omega_2
&=\int_0^1 0\times1\,\mathrm{d}t+\int_0^1 t\times0\,\mathrm{d}t=0.
\end{aligned}
\]
$\gamma_2$ के पहले खंड पर $\mathrm{d}x=0$ और दूसरे पर $x(t)=t$, $y(t)=1$ है, इसलिए
\[
\begin{aligned}
J_2=\int_{\gamma_2}\omega_2
&=\int_0^1 t\times0\,\mathrm{d}t+\int_0^1 1\times1\,\mathrm{d}t=1.
\end{aligned}
\]
विकर्ण खंड $\gamma_3$ के लिए
\[
x(t)=t,\qquad y(t)=t,\qquad t\in[0,1],
\]
अतः
\[
J_3=\int_{\gamma_3}\omega_2
=\int_0^1y(t)x'(t)\,\mathrm{d}t
=\int_0^1t\,\mathrm{d}t=\frac12.
\]
समान सिरों के बावजूद $J_1=0$, $J_2=1$, $J_3=1/2$ हैं। यही गैर-यथार्थ अवकल रूप की पथ-निर्भरता है।
\end{solution}
\end{exo}

\begin{exo}[आंतरिक ऊर्जा का समान परिवर्तन, अंतरण के तीन ढंग]{meme-delta-u-trois-transferts}
\keywords{ऊष्मागतिकी का प्रथम नियम, आंतरिक ऊर्जा, ऊष्मा, यांत्रिक कार्य, विद्युत कार्य, ऊष्मामिति, अवस्था फलन}

एक द्रव, उसका ऊष्मामापी और उसमें डूबा छोटा प्रतिरोधक मिलकर एक बंद, स्थूलतः विरामस्थ निकाय बनाते हैं।
पात्र दृढ़ है और निकाय की कुल ऊष्मा धारिता अचर मानी गई है:
\[
C=2{,}00\ \mathrm{kJ\,K^{-1}}.
\]
निकाय को $T_A=20{,}0\,^\circ\mathrm{C}$ वाली साम्यावस्था $A$ से
$T_B=25{,}0\,^\circ\mathrm{C}$ वाली साम्यावस्था $B$ में ले जाना है। मानें कि
\[
U(B)-U(A)=C(T_B-T_A).
\]
स्थूल गतिज और स्थितिज ऊर्जाओं के परिवर्तन शून्य हैं तथा बाहरी हानियाँ उपेक्षित हैं।

उसी प्रारंभिक अवस्था $A$ से निम्न तीन विधियाँ स्वतंत्र रूप से की जाती हैं:
\begin{itemize}
\item[विधि 1] कोई कार्य दिए बिना निकाय को गर्म स्रोत के ऊष्मीय संपर्क में रखा जाता है।
\item[विधि 2] दीवारें ऊष्मारोधी हैं। $m$ द्रव्यमान का भार $h=5{,}00$ m गिरकर पैडल से द्रव को हिलाता है। अंत में पैडल और द्रव विराम में हैं तथा भार की खोई सारी यांत्रिक ऊर्जा निकाय को मिलती है। $g=9{,}81\ \mathrm{m\,s^{-2}}$ लें।
\item[विधि 3] निकाय को ऊष्मा $Q_3=4{,}00$ kJ मिलती है और शेष ऊर्जा प्रतिरोधक को विद्युत रूप से दी जाती है। प्राप्त विद्युत शक्ति अचर $\mathcal P=300$ W है।
\end{itemize}

\begin{enumerate}
\item तीनों विधियों के लिए समान आंतरिक ऊर्जा परिवर्तन $\Delta U=U(B)-U(A)$ निकालें।
\item पहली दो विधियों में निकाय को मिली ऊष्मा $Q$ और कार्य $W$ निकालें; दूसरी में आवश्यक द्रव्यमान $m$ भी निकालें।
\item तीसरी विधि में प्राप्त विद्युत कार्य और प्रतिरोधक को शक्ति देने की अवधि निकालें।
\item आरंभिक और अंतिम अवस्थाएँ समान होने पर भी $Q$ और $W$ अलग क्यों हो सकते हैं? क्या अंतिम अवस्था $B$ में निकाय के भीतर “ऊष्मा” या “कार्य” होता है?
\end{enumerate}

\begin{indication}
बैंकर चिह्न परंपरा से $\Delta U=Q+W$ लगाएँ। पैडल में प्राप्त कार्य भार की स्थितिज ऊर्जा की कमी $mgh$ है। तारों से सीमा पार करती विद्युत ऊर्जा को विद्युत कार्य गिनें।
\end{indication}

\begin{solution}
\textbf{प्रश्न 1।} $T_B-T_A=5{,}0$ K, अतः
\[
\Delta U=C(T_B-T_A)=2{,}00\times5{,}0=10{,}0\ \mathrm{kJ}.
\]
यह केवल साम्यावस्थाओं $A$ और $B$ पर निर्भर करता है।

\textbf{प्रश्न 2।} विधि 1 में $W_1=0$, इसलिए प्रथम नियम से
\[
Q_1=\Delta U=10{,}0\ \mathrm{kJ}.
\]
ऊर्जा निकाय में आने से ऊष्मा धनात्मक है। विधि 2 में दीवारें ऊष्मारोधी हैं, अतः $Q_2=0$ और
\[
W_2=\Delta U=10{,}0\ \mathrm{kJ}.
\]
$W_2=mgh$ से
\[
m=\frac{W_2}{gh}
=\frac{10{,}0\times10^3}{9{,}81\times5{,}00}
\simeq204\ \mathrm{kg}.
\]
यह बड़ा परिमाण बताता है कि मामूली तापवृद्धि भी पर्याप्त यांत्रिक ऊर्जा के तुल्य है।

\textbf{प्रश्न 3।} प्रथम नियम से
\[
W_3=\Delta U-Q_3=10{,}0-4{,}00=6{,}00\ \mathrm{kJ}.
\]
विद्युत आपूर्ति निकाय को ऊर्जा देती है, इसलिए कार्य धनात्मक है। $W_3=\mathcal P\,\Delta t$ से
\[
\Delta t=\frac{W_3}{\mathcal P}=\frac{6{,}00\times10^3}{300}=20{,}0\ \mathrm{s}.
\]

\textbf{प्रश्न 4।} तीन पथ क्रमशः
\[
(Q_1,W_1)=(10{,}0,0)\ \mathrm{kJ},\qquad
(Q_2,W_2)=(0,10{,}0)\ \mathrm{kJ},
\]
और
\[
(Q_3,W_3)=(4{,}00,6{,}00)\ \mathrm{kJ}
\]
देते हैं। हर बार $Q+W=10{,}0$ kJ और $\Delta U$ समान है। आंतरिक ऊर्जा अवस्था फलन है,
जबकि ऊष्मा और कार्य प्रक्रम में ऊर्जा-अंतरण बताते हैं। अंतिम अवस्था $B$ में निकाय की आंतरिक ऊर्जा है,
पर उसमें न “ऊष्मा” है न “कार्य”।
\end{solution}
\end{exo}

\begin{exo}[अचर बाह्य दाब में संपीडन]{compression-pression-exterieure-constante}
\seoready{true}
\keywords{दाब बलों का कार्य, अचर बाह्य दाब, संपीडन, प्रसार, निर्वात में प्रसार}

गैस को अचर बाह्य दाब $P_0$ में आयतन $V_A$ से $V_B<V_A$ तक संपीडित किया जाता है।

\begin{enumerate}
\item गैस द्वारा प्राप्त कार्य निकालें।
\item गैस उसी बाह्य दाब $P_0$ के विरुद्ध $V_B$ से $V_A$ लौटती है। प्राप्त कार्य निकालकर संपीडन के कार्य से तुलना करें।
\item निर्वात में $V_B$ से $V_A$ तक प्रसार को दोहराएँ और तीनों चिह्नों की व्याख्या करें।
\end{enumerate}

\begin{indication}
$W=-\int P_{\mathrm{ext}}\,\mathrm{d}V$ का प्रयोग करें। प्रक्रम अर्ध-स्थैतिक न हो तब भी \emph{बाह्य} दाब का ही समाकलन करना है।
\end{indication}

\begin{solution}
\textbf{प्रश्न 1।} बाह्य दाब अचर है, अतः
\[
W_{A\to B}=-P_0(V_B-V_A)=P_0(V_A-V_B)>0.
\]
संपीडन में गैस कार्य प्राप्त करती है।

\textbf{प्रश्न 2।} वापसी में
\[
W_{B\to A}=-P_0(V_A-V_B)<0.
\]
दोनों कार्य विपरीत हैं क्योंकि प्रक्रम समान अचर बाह्य दाब के विरुद्ध होते हैं।

\textbf{प्रश्न 3।} निर्वात में $P_{\mathrm{ext}}=0$, अतः
\[
W_{B\to A}^{\mathrm{vide}}=0.
\]
प्रसार अपने-आप ऋणात्मक कार्य नहीं देता; वास्तव में किसी बाह्य बल को पीछे धकेलना आवश्यक है।
\end{solution}
\end{exo}
